\scanpage{48}
Cela, plus le fait que $t_G$ et $\Omega^{\vee}\otimes_{\mathbb Z_p}A$ sont de rang $1$ sur $A$ resp. $A\otimes_{\mathbb Z_p}A$, et que $\unsure{\rho|\varphi}$ est surjectif, prouve que $t_G$ \uline{est déduit de} $\Omega^{\vee}\otimes_{\mathbb Z_p}A$ par le changt. de base
\diagram{\begin{tikzcd}A\otimes_{\mathbb Z_p}A\arrow[r,no head,"\rho_0"]&A\end{tikzcd}}
donc \deleted{\unsure{par suite}} (comme $\Omega^{\vee}\otimes_{\mathbb Z_p}A$ est déduit de $\Omega^{\vee}$ par $A\xrightarrow{\lambda\longrightarrow\lambda\otimes1}A\otimes_{\mathbb Z_p}A$) \unsure{on} obtient un isom can.
\[
t_G\simeq\Omega^{\vee}\qquad\text{i.e.}\qquad\omega_G\simeq\Omega
\]
et la suite exacte $(**)$ est la suite exacte canonique évidente, \inserted{d'où $(*)$ par dualité.} Comme cette filtration étend celle de $M_A\otimes_A k$, on a retrouvé ici le th. d'existence et d'unicité de Lubin-Tate, (sous une forme moins forte d'ailleurs en ce qui concerne l'unicité, et sous l'hyp. complémentaire que $\exists$ puiss. div. sur $\pi A$.

Si on abandonne cette dernière hypothèse, et qu'on utilise la classification des déformations à isogénie près, on trouve de $\unsure{=}$ qu'on a une suite exacte canonique
\diagram{\begin{tikzcd}
0\arrow[r,no head]&\omega_G\otimes_A K\arrow[r,no head]&M_K\arrow[r]\arrow[d,equal]&t_{G^*}\otimes_{\unsure{A}}K\arrow[r,no head]&0\\
&&\Omega\otimes_{\mathbb Z_p}K
\end{tikzcd}}
ou encore
\diagram{\begin{tikzcd}
0\arrow[r,no head]&\omega_{G^*}\otimes_A K\arrow[r,no head]&M_K^{\vee}\arrow[r]\arrow[d,equal]&t_G\otimes_A K\arrow[r]&0\\
&&\Omega^{\vee}\otimes_{\mathbb Z_p}K
\end{tikzcd}}
et on trouve encore que
\[
t_G\otimes_A K\simeq M_K^{\unsure{\vee}}\otimes_{\unsure{\Omega^{\vee}\otimes_{\mathbb Z_p}K}}(K\otimes K,\rho_0)\simeq\Omega^{\vee}\otimes_A K
\]
\mspage{47}
